(1) 定态方程
\[
-\frac{\hbar^2}{2m}\psi''(x)-g\delta(x)\psi(x)=E\psi(x)
\]
在原点以外是自由粒子方程。波函数在有限 \(\delta \) 势下连续，\(\psi(0^+)=\psi(0^-)=\psi(0)\)；对 \([-\varepsilon,\varepsilon]\) 积分并令 \(\varepsilon\downarrow0\)，得到导数跃变
\[
\psi'(0^+)-\psi'(0^-)=-\frac{2mg}{\hbar^2}\psi(0)=:-2\kappa\psi(0),
\quad\kappa=\frac{mg}{\hbar^2}>0.
\]

(2) 若束缚态 \(E=-\hbar^2q^2/(2m)<0\)，归一化要求 \(\psi(x)\) 在两个无穷远均衰减，因此有 \(\psi(x)=A e^{-q|x|}\)。导数跃变为 \(-2qA\)，与边界条件对比得 \(q=\kappa\)，故
\[
E_b=-\frac{\hbar^2\kappa^2}{2m}=-\frac{mg^2}{2\hbar^2},\qquad
\psi_b(x)=\sqrt\kappa\,e^{-\kappa|x|}.
\]
归一化用 \(\int_{\mathbb R}e^{-2\kappa|x|}dx=1/\kappa\)。奇宇称函数在 \(x=0\) 为零，跳跃条件退化为导数连续，而左右无穷远同时衰减的非零奇态不可能满足此条件。因此恰有一个偶束缚态。

(3) 对 \(E=\hbar^2k^2/(2m)>0\)，写
\[
\psi(x)=
\begin{cases}
e^{ikx}+r e^{-ikx},&x<0,\\
t e^{ikx},&x>0.
\end{cases}
\]
连续性给 \(t=1+r\)，导数跳跃给 \(2ikr=-2\kappa(1+r)\)，故
\[
r=-\frac{\kappa}{\kappa+ik},\qquad
t=\frac{ik}{\kappa+ik}.
\]
两侧自由粒子的波数相同，概率流与振幅模平方成正比，因此
\[
R=|r|^2=\frac{\kappa^2}{k^2+\kappa^2},\qquad
T=|t|^2=\frac{k^2}{k^2+\kappa^2},\quad R+T=1.
\]

(4) 当 \(g\to0^+\)，\(\kappa\to0\)，束缚态能量趋 \(0^-\) 且空间衰减长度 \(1/\kappa\to\infty\)；对固定 \(k>0\)，散射 \(R\to0\)。当 \(g\to\infty\)，\(\kappa\to\infty\)，对固定能量 \(R\to1\)。对任何固定 \(g>0\)，低能 \(k\to0^+\) 同样有 \(R\to1\)，说明无限长波对任意弱但非零的吸引 \(\delta \) 势也会近似全反射；\(k\to0\) 与 \(g\to0\) 的极限次序不能随意交换。
